\chapter[Sermon 80. The Rejoicing of Saints for the Overthrow of Babylon, the Drowning of the Same, and the Causes of Drowning or Destruction Are Rehearsed.]{The Rejoicing of Saints for the Overthrow of Babylon, the Drowning of the Same, and the Causes of Drowning or Destruction Are Rehearsed.}
\label{sermon:080}
\markright{Sermon 80. The Rejoicing of Saints for the Overthrow of Babylon ...}

Rejoice over her, heaven, and the holy Apostles and prophets, for God has given you judgment upon her. And a mighty angel lifted up a great millstone and cast it into the sea, saying, “With such violence will that great city Babylon be cast down, and will be found no more.” And the sound of harps and musicians, and of pipers and trumpets, will be heard no more in you; and no craftsman, whatever his craft may be, will be found no more in you; and the sound of a mill will be heard no more in you; and the light of a candle will shine no more in you; and the voice of the bridegroom and of the bride will be heard no more in you, for your merchants were princes of the earth, and with your enchantments all nations were deceived; and in her was found the blood of the prophets and of the saints, and of all those who were slain upon the earth.

In the fifteenth place of this chapter, the Angel of the Lord exhorts all the saints in Heaven to rejoice, and that for the overthrow of Babylon. And this rejoicing of the saints is set against the wailing of the wicked. For as they lament for the loss of pleasures, so the saints rejoice over ungodliness oppressed, and the glory of God avenged.

We are certainly forbidden in the Proverbs of Solomon, and in the doctrine of Christ and his Apostles, that we should not be glad of the calamities of our enemies, nor that we should say evil or do evil to our enemies. This is a perpetual command, given to all men, never to be altered by any dispensation. But we must observe that men rejoice in diverse ways. Men are glad many times at the destruction of their enemies, and that out of hatred and malice, which is not done without sin. Others are glad again of the calamities and plagues of the ungodly, yet bearing no malice toward them, and would doubtless have wished them a better state if they might have been persuaded to turn. But they rejoice rather over justice revenged, and the godly delivered from the tyranny of the wicked. Of this we read that the prophet said in Psalm 58: “The righteous shall rejoice when he shall see vengeance; he shall wash his feet in the blood of the ungodly” (that is, he shall purge his affections and evil manners, when he shall see the blood of the ungodly spilled, which he believes to be done as a warning, lest we should follow our evil affections, and that our blood should be shed also by the most just God through his ministers). And a man will say, “Verily, there is a reward for the righteous; verily, God judges the earth.” Therefore, the righteous are glad and rejoice when they see vengeance. And it is not said that they covet or wish for vengeance. “Vengeance is mine,” says the Lord, “I will reward.” When the Lord therefore rewards, they are glad for the deliverance, and for the truth established and confirmed, and rejoice not out of hatred they have toward the oppressors, whom they have wished lost and destroyed.

\index{Repentance}The godly always wish the wicked to be converted and to return into favor with God. But when they see them unmoved by repentance, but obstinately proceeding and falling into their own destruction, and that God intervenes for the salvation of the faithful and the deliverance of the godly, the godly rejoice at this deliverance and praise the justice of God. Nevertheless, they would always rather, if it could be, that the lost had led their lives otherwise. But now, since it cannot be otherwise through their own obstinate malice, they do not speak against the judgments of God, but rather commend them. These things, certainly, the saints on Earth do. And the saints in Heaven, since they are now purified from all affections, their rejoicing is altogether most pure, so that it would be superfluous to reason curiously about it.

But where the heavenly saints rejoice at the destruction of the wicked, we may easily judge how much they err who trust to the help or prayers of saints, and yet alter nothing at all of their wicked life. It shall be easy also to discuss their doubt and carefulness, who fear lest their brethren, sisters, friends, and kinsfolk be condemned. For the saints plainly consent to the will of God, and extol the judgments of God, and rejoice thereat, and can be sorry no more.

And he bids them rejoice, as many times in the Psalms we read a like phrase; unless you would rather understand by “Heaven,” heavenly dwellers, such as we believe the apostles and prophets to be. For at the same time when John wrote these things, all the apostles were in a manner slain. And here it is to be known that the Roman beast had devoured, that is to say, afflicted and slain, not only the Son of God, our Lord Jesus Christ, but also John the Baptist, all the Apostles of God, and all the martyrs of Christ. By the prophets we understand, not only those of old, but all the faithful preachers of the gospel. For we have heard more than once that faithful preachers of the word are called prophets. He adds moreover a reason why they ought to rejoice: for God has given you judgment concerning her. For in chapter 6, the souls of the martyrs cry under the altar: “Lord, how long before you avenge our blood on those who dwell on Earth?” Now therefore they praise God’s justice, which, as he promised he would avenge, so has he now avenged in deed.

\index{Rome}And by this passage we learn that all judgment is given to the Son, and that no saint in heaven can judge or punish an evil man on Earth. For it is most false that saints are said to punish their enemies: Saint Anthony with the holy fire, Valentine with the falling sickness, and others with other diseases; God alone, as is declared at large in chapter 16, punishes and sends and takes away sickness. And most certain it is, both by this and by many other passages in this book, that God does not sleep, but will, when he sees fit, certainly avenge and punish. The martyrs, when they were about to die, had committed all their judgment to the Lord their God. He now judges the judgment of the saints of Rome: that is, after his just judgment takes punishment of Rome, for that she had with wrongful judgment oppressed the saints.

\index{Primasius}\index{Papacy}In the sixth place of this chapter he returns to the description of the subversion of Babylon. And it is a most clear, and even a certain visible and evident demonstration by a similitude. For taking up a great stone, in quantity like a millstone, he casts the same into the sea, and making a declaration of his doing, says, thus suddenly, and with such a violence, shall Babylon be cast down, \&c. This place is taken out of the end of the 51st chapter of Jeremiah, where you read in a manner the like things word for word. And here is now brought in a strong angel, lest we should think that the force of Rome were happily stronger than that it could be broken. But it shall be broken by a strong angel. And the things that be suddenly drowned appear no more. Here is signified therefore, that with a sudden destruction Rome shall fall, that there shall no token thereof be left, and that it shall fall without any difficulty, it shall be made to plunge, and never more be seen. And the Lord in the gospel affirms, that the crime or slander must be punished with a millstone hanged about the neck: yea and that same not to be punishment grievous enough, although among the Syrians it was accounted for vile and shameful, since the crime deserves to be punished with a much more grievous or crueler pain. Wherefore Primasius supposed, that here by the way is signified how Babylon, for offenses given to the world, should be drowned in the sea, as it were with a millstone tied fast to her neck. Doubtless if ever any city, if ever any kingdom were hateful by reason of greatest offenses, and given to the Christians innumerable slanders: Rome and the Roman Empire, and even the papists of the church have hurt most by slander, and yet hurt. Wherefore it is no doubt, but that it has been plagued most grievously, and shall be yet more punished by the Lord.

\index{Ezekiel (Book of)}Again, through prophetic and figurative language, he signifies a significant desolation, and that the place would never afterward be inhabited forever. You will find similar expressions in Isaiah 24 and Ezekiel 26, and in various other passages. All pleasure, he says, will perish, especially that which was formerly derived from music. All trades will cease. Briefly, there will be no more any dwelling place for people.

\index{Arethas of Caesarea}In the seventh and final place, I set forth again the causes of this downfall, and these are particularly noteworthy. The first: “Your merchants were princes of the earth.” For those who have engaged in commerce within the church of Rome, and still do, are in a manner princes. Of whom I have spoken before. Here is noted therefore their pride, avarice, and lavishness. Arethas says he calls them merchants, who trouble and vex the entire world, as if they were certain fairs, and so forth. The second: “By your enchantment all nations have been seduced.” There is no doubt that enchantment and magic reign in Babylon, and that there are plenty of fortune-tellers, necromancers, and enchanters to be found there; yet here chiefly is signified seduction, idolatry, and impiety, or error of doctrine. Such an enchanter was Jezebel, as appears in the fourth book of Kings, chapter 9, who practiced enchantment indeed, and bewitched people with corrupt religion. And even so Rome has seduced the entire world, and yet seduces. For this cause she deserves most grievous punishment. The last cause of subversion: “For in you is found the blood.” Blood shed cannot be wiped away nor cleansed from those who shed innocent blood. And although it is not immediately required, there will come a time when it shall be required of God, and then it will be found. And he makes mention of three kinds of blood. First, the blood of prophets, meaning those who have preached the Gospel and have been the fathers of the faithful. Secondly, the blood of saints, namely, holy martyrs. Finally, the blood of all men who have been slain on earth, namely, dwelling here and there throughout the world, whom we understand to have been dispatched and taken out of the way by the wars, seditions, and tyranny of Rome. So we also read in the first oration of Jeremiah, that God straightly requires the blood of his servants spilled. Doubtless all shedding of blood is grievous (the same excepted which is justly done by the magistrate), yet one is more heinous than another. For he who kills a preacher of the Gospel sins more grievously than he who dispatches a private person, and he who for religion’s sake slays a man, making him a martyr, sins more heinously than he who kills a man in war. Therefore, all the blood shed by Rome in any sort shall be required of Rome, and is required. Thus the Lord also spoke of the city of Jerusalem, Matthew 23. The Lord Jesus have mercy on us, and look upon us with eyes of his mercy. Amen.
